\begin{abstract}
    \noindent
High-quality LLM request scheduling requires
achieving two key objectives: whether the routed instance has
{\kvcache} to accelerate the request execution and whether
the workload is balanced across instances.
Achieving both objectives is challenging because
pursuing one objective may compromise the other.
Current approaches adopt various combinators (e.g., linear combinations)
to compute a scheduling score combining indicators for the two objectives,
which are complex in that they either require significant workload-specific hyperparameter tuning
or model-hardware-aware simulator development,
and could still lead to suboptimal performance.

In this paper, we show that using a simple multiplication of two carefully chosen indicators---one 
for {\kvcache}-aware (new prefill tokens if routed to an instance) and one for load balancing-aware 
(current batch size of the instance)---as the scheduling score
can simultaneously achieve both objectives well without any hyperparameter tuning.
The key idea is that the multiplied score considers both objectives in a manner similar to a linear combination,
with a nice property that the original hyperparameters are canceled out during comparison 
so we don't need tuning to find the best parameters.  
The two indicators are chosen based on our analysis of LLM characteristics, 
and our extensive experiments show that this simple approach can reduce TTFT by 92\% and 52\%, and TPOT by 21\% and 20\%, 
compared to vLLM-v1 and a production scheduler on real-world workloads covering chatbots, API calls, 
and coding agents.
We also mathematically derive the conditions under which multiplication may fail, 
and find that such conditions are extremely rare in practice and 
can be detected (and mitigated) beforehand.

\end{abstract}
